\chapter{Theoretical Foundations and Literature Review}
\label{ch:foundations}

This chapter sets out the bodies of theory the thesis relies on --- latency
arbitrage and discrete-time market design, adverse selection, market
transparency, uniform-price auction theory and the measurement of market
quality --- and closes with the price-improvement argument that motivates
the boundary-timing problem of batch auctions.

\section{Discrete-Time Market Design}
\label{sec:discretetime}

\subsection{The Mechanics of Latency Arbitrage}

The argument of \citet{budish2015hft} rests on a single structural
observation, not on irrationality or market power: in a continuous
market, a public information event creates a discrete arbitrage
opportunity that is allocated by relative speed.

Consider a liquidity provider quoting a two-sided market in an asset
whose fair value is tracked by a correlated, continuously observable
signal. When the signal jumps, every participant learns the same thing at
essentially the same instant: the liquidity provider wants to cancel its
stale quote, every other fast participant wants to hit it. The engine
processes both messages in arrival order, so whether the quote is
cancelled or picked off depends on network path length, hardware and code
efficiency, not on anything informational.

Two consequences follow. The opportunity is mechanically recurrent,
appearing every time the public signal moves --- thousands of times a day
in a liquid instrument. And the value of being marginally faster does not
decline as everybody invests in speed: since only the winner is paid, the
equilibrium is a rent-dissipating race in which expenditure on latency is
bounded by the value of the opportunities, which persist however fast the
market becomes. Speed is a purely relative good.

The cost is not borne by the racers. A liquidity provider that expects to
lose a fraction of its volume to picked-off quotes must recover that loss
from the trades it does not lose, shifting the burden onto liquidity
demanders as a wider spread --- a tax on liquidity in the words of
\citet{budish2015hft}.

\subsection{Empirical Magnitude}

\citet{aquilina2022quantifying} measure this directly, using London Stock
Exchange message-level data (including messages never executed) to
identify latency races as clusters of near-simultaneous attempts to take
or cancel the same quote. Races are extremely frequent and short --- far
shorter than any batch interval measured in milliseconds; a large share of
the volume at the exact moment of a price change is race volume; and the
aggregate transfer, small relative to total volume but large in absolute
terms, concentrates in the effective spread paid by liquidity demanders.
Removing the race therefore removes a component of the spread.

\subsection{Batching as the Response}

Frequent batch auctions change two rules at once, and the two do
different work. \emph{Discrete time} removes the informational value of
within-interval arrival order: if orders are collected over $[t,
t+\tau)$ and matched only at $t+\tau$, any two orders inside the window
are symmetric with respect to time. \emph{Uniform pricing} removes the
value of queue position within the matched set: if all executions occur
at one price $\clr$, no participant obtains a better price by being
earlier. Either rule alone is insufficient --- discrete time with
discriminatory pricing would still reward earlier submission, and uniform
pricing in continuous time is not well defined. Together they imply that
speed pays only for crossing an interval boundary, getting into batch $n$
rather than batch $n+1$: the value of a marginal nanosecond falls from
``wins the race'' to ``matters only in the vanishingly rare case that it
changes which batch an order lands in''. In the model of
\citet{budish2015hft}, liquidity providers then no longer need to price
sniping risk into their quotes and spreads narrow. The residual cost is
delay, an average of $\tau/2$ if orders arrive uniformly, so the choice of
$\tau$ trades speed neutralisation against delay imposed on everybody.

\section{Information Asymmetry and Adverse Selection}
\label{sec:adverse}

Latency arbitrage is one specific mechanism through which a liquidity
provider loses money to better-informed counterparties; the general
theory of that loss predates high-frequency trading by three decades and
is what makes the spread a meaningful economic object.

\subsection{The Spread as an Information Premium}

\citet{glosten1985bid} model a competitive, risk-neutral market maker
facing a mixed population of informed and uninformed traders it cannot
tell apart. Because a buy order is more likely to come from an informed
trader when the asset is undervalued, every order carries information, so
quotes must be conditionally fair given the direction of the order
expected --- a bid--ask spread that exists purely as an information
premium, even with no inventory risk or processing cost. The spread widens
with the proportion of informed traders and the precision of their
information, and narrows when the market maker can identify a subset of
flow as uninformed. This is the theoretical foundation of order flow
segregation: separating benign flow lets a liquidity provider quote
tighter against it without cross-subsidising losses on the rest.

\subsection{Strategic Informed Trading and Market Depth}

\citet{kyle1985continuous} approaches the same problem from the informed
trader's side: a single informed agent chooses how much to trade, knowing
that trading more reveals more, while noise traders provide cover and a
competitive market maker sets the price as the conditional expectation of
value given total order flow. The equilibrium price function is linear,
$p = \mu + \lambda \cdot y$, where $y$ is aggregate order flow and
$\lambda$ measures price impact per unit of flow; its reciprocal, market
depth, is the volume required to move the price by one unit. Kyle's model
is a batch model --- the market maker observes aggregate flow and sets a
single price, the structure of a uniform-price call auction --- so the FBA
literature can be read as asking what happens when a real market is
organised the way Kyle already assumes. $\lambda$ also supplies the
empirical bridge: price impact is the operational measure of adverse
selection.

\subsection{Decomposing the Spread}
\label{sec:spread-decomposition}

The adverse selection premium is not directly observable; only the total
cost of a trade is. The standard strategy is to decompose the effective
spread into a realised spread and a price impact component
\citep{hasbrouck1993assessing,hendershott2011algorithmic}. Let $p_k$ be
the execution price of trade $k$, $D_k \in \{+1,-1\}$ its direction
(buyer- or seller-initiated), $m_k$ the midpoint immediately before the
trade, and $m_{k+\Delta}$ the midpoint a markout horizon $\Delta$
afterwards. In proportional form,
\begin{align}
  S^{\mathrm{eff}}_k &= \frac{2 D_k (p_k - m_k)}{m_k},
    \label{eq:effspread}\\[3pt]
  S^{\mathrm{real}}_k(\Delta) &= \frac{2 D_k (p_k - m_{k+\Delta})}{m_k},
    \label{eq:realspread}\\[3pt]
  \mathrm{PI}_k(\Delta) &= \frac{2 D_k (m_{k+\Delta} - m_k)}{m_k},
    \label{eq:priceimpact}
\end{align}
so that $S^{\mathrm{eff}}_k = S^{\mathrm{real}}_k(\Delta) +
\mathrm{PI}_k(\Delta)$ trade by trade and in any volume-weighted
aggregate. The effective spread is what the liquidity demander paid relative to the
midpoint; the realised spread is what the provider still holds once the
market has revised its valuation, the gross compensation for supplying
liquidity; the price impact is the permanent revision of the midpoint in the
direction of the trade, the part lost to the information the order carried.

When a public signal moves, discrete uniform-price clearing removes the
latency-arbitrage mechanism, since the aggressive order and the maker's
cancellation are processed together and execution occurs at a price that
already reflects orders submitted in response to the signal. It does not
remove adverse selection in the theoretical sense --- a genuinely
informed trader can still price aggressively into a batch
\citep{eibelshauser2022fba} --- so the prediction is a reduction in
price impact, not its elimination. A fall in effective spread driven by
price impact is therefore the result theory predicts; a fall driven by
realised spread would mean providers simply earn less with the
informational component untouched, and a fall in price impact with no
change in effective spread would mean the benefit is captured by
providers rather than passed on.

\section{Market Transparency Theory}
\label{sec:transparency}

\citet{madhavan1996transparency} defines transparency as participants'
ability to observe information about the trading process, and shows more
of it is not unambiguously better: greater pre-trade transparency helps
prices aggregate information and reduces uncertainty about market state,
but also exposes the trading intent of participants transacting in size,
and can increase volatility and reduce liquidity in some configurations.
Consolidating with \citet{ohara1995microstructure}, disclosure design is
a genuine trade-off, not a monotone improvement.

In an FBA the question is narrower: what is disclosed during the
accumulation phase, before the auction clears? Moving toward opacity
protects submitted orders from being exploited but degrades the
information available for pricing, and classical auction theory predicts
that participants respond to uncertainty about the clearing price by
shading their bids \citep{klemperer1999auction}. In a uniform-price
setting shading is not merely a transfer: shaded limit prices are less
likely to cross, so it can reduce the volume that clears.

The trade-off assumes the operator can choose where to sit among these
options. Decentralised blockchain venues are the limiting case where it
cannot: a transaction waiting in a public mempool is visible to everyone
before it executes, so such markets are fully open by default without the
benefit that regime is supposed to buy. Disclosure normally exposes the trading intent of the order submitter
\citep{madhavan1996transparency}, and the informational disadvantage this
creates normally falls on the market maker posting a quote
\citep{glosten1985bid}. A mempool flips both of these. What is exposed is
not a resting quote but an instruction that has not executed yet, and the
party exposed by it is not the quoter but the one demanding liquidity,
the submitter itself. An observer can buy ahead of a pending swap and
sell into the price it moves, sandwiching the submitter. Batch
clearing removes this by construction: if every order in the interval
settles at one uniform price, position within the batch carries no value.

\section{Auction Theory and Uniform-Price Clearing}
\label{sec:auctiontheory}

In a uniform-price auction all winning participants transact at the same
market-clearing price, independent of the prices they submitted. A
participant's own bid determines \emph{whether} it trades but, except for
the marginal order, not \emph{at what price}. Misreporting has no upside:
bidding below one's valuation only risks losing an acceptable trade,
bidding above only risks winning an unacceptable one, so for a
price-taking participant with single-unit demand truthful revelation is
approximately optimal, a property inherited from
\citet{vickrey1961counterspeculation}. This is the formal basis for the
claim that batching substitutes competition in price for competition in
speed: if submitted limits are approximately truthful, the clearing price
aggregates valuations rather than reaction times. Visibility of the
accumulated batch does not disturb this --- a participant who can
anticipate the clearing price pays it regardless of the limit it submits
above it --- but it does change \emph{when} orders are submitted, since
observing the indicative price gives a participant reason to withhold
until shortly before the batch closes.

The qualification carried by ``approximately'' depends jointly on
pivotality and size. A participant too small to move the crossing point
controls only its own inclusion, so shading saves nothing and truthful
bidding is optimal without qualification. A pivotal participant with
single-unit demand does pin the clearing price with its own limit, but the
marginal unit carries almost no surplus, so the incentive to shade is
negligible. The case that matters is the pivotal participant with
multi-unit demand: its last unit sets the price for every unit it buys, so
shading the marginal bid reduces the cost of the entire inframarginal
block while risking only the one unit barely worth trading. The saving
scales with size, the loss does not, and the incentive to misreport is
restored; \citet{klemperer1999auction} surveys this literature. Pivotality
is relative to the liquidity accumulated within one batch, not the market
as a whole, so the strength of the truthfulness property is itself a
function of interval length: shorter intervals produce thinner batches,
where an order that would be atomistic in a long window can be
price-setting in a short one.

For every executed non-marginal order, uniform pricing also produces price
improvement relative to the submitted limit: a buyer willing to pay
$\pi_k$ who clears at $\clr < \pi_k$ captures $(\pi_k - \clr) q_k$ of
surplus that, under continuous discriminatory execution at the resting
price, would in general have accrued differently. Aggregating this across
executed orders gives the realised trader surplus metric.

\section{Measurement of Market Quality}
\label{sec:quality}

The empirical comparison relies on the standard microstructure toolkit.
Quoted spread and depth describe the cost a hypothetical trade would face;
effective spread describes what an actual trade cost; realised spread and
price impact decompose that cost into the part retained by the liquidity
provider and the part lost to adverse selection
\citep{hasbrouck1993assessing}. \citet{amihud2002illiquidity} provides a
low-frequency summary of price impact per unit of volume, robust and easy
to compare across mechanisms, and \citet{hasbrouck1995onesecurity}
supplies the framework for attributing price discovery across venues.
Hyperliquid publishes its own oracle price for margin and liquidation, but
the order-status dataset replayed here does not carry it, so pricing error
against an external reference is defined but uncomputable in this thesis.

Two adaptations are required for a batch mechanism. Measures that reference
the prevailing midpoint at the moment of the trade have no book to take a
midpoint from during accumulation, so the reference is the previous batch's
own clearing price as of the order's submission --- stale by up to $\tau$,
never a same-instant quantity as the CDA's midpoint is. And the quoted
spread has no direct analogue during accumulation, defined here as the
implied spread between the marginal unexecuted buy and sell of the batch.

Finally, the experimental convention is that of \citet{gode1993allocative}:
hold the induced order flow constant and vary only the mechanism. Their
demonstration that institutions can produce allocative efficiency largely
independently of trader rationality also licenses replay-based simulation,
since much of what a mechanism does is mechanical rather than behavioural;
\citet{lebaron2006agent} surveys the agent-based alternative in which
behaviour is modelled and allowed to adapt.

\newpage
\section{Price Improvement produces the Boundary
Repricing Race}
\label{sec:priceimprovement}

The argument developed here is, to the author's knowledge, not stated in
this form in the existing literature. Its components are individually
familiar: the residual value of low latency near a batch boundary is
established by \citet{schlegel2023transaction}, the tendency of
sophisticated traders to submit late in transparent call auctions is
documented by \citet{daures2026click}, and the strategic consequences of
the clearing-price selection rule in uniform-price auctions are studied by
\citet{burkett2020uniform}. What has not been made is the connection
between them. The frequent batch auction literature treats the price
improvement that uniform clearing confers on inframarginal orders as a
benefit of the design; this section argues that the treatment rests on a
category error. Price improvement measured on one side of the book is not
a welfare quantity but a transfer, and the rule that determines its size
is manipulable by the timing of order submission. The FBA therefore does
not remove the incentive to compete on speed: it converts a race for
priority into a race for proximity to the batch boundary, independently of
the informational motives usually invoked to explain late submission.

\subsection{Setup}

This argument considers a single batch, using the demand and supply
schedules $D(p)$/$S(p)$, the volume-maximising set $\mathcal{P}^{*}$, and
the clearing price $\clr = \sigma(\mathcal{P}^{*})$ defined in
\Cref{sec:fbanotation}, plus the executed quantity $x_k \ge 0$ of order
$k$ under a volume-maximising allocation, so that $\sum_{k \in \mathbb{B}}
x_k = \sum_{k \in \mathbb{A}} x_k = V^{*}$.

\subsection{Price Improvement Is a Transfer}

Write the realised surplus of the two sides at clearing price $p$ as
\begin{equation}
  \mathrm{PI}^{\mathbb{B}}(p)
    = \sum_{k \in \mathbb{B}} (\pi_k - p)\, x_k,
  \qquad
  \mathrm{PI}^{\mathbb{A}}(p)
    = \sum_{k \in \mathbb{A}} (p - \pi_k)\, x_k .
\end{equation}

\textbf{Total surplus does not depend on which clearing price is chosen.}
For any $p \in \mathcal{P}^{*}$ with allocation $x$ held fixed,
\begin{equation}
\label{lem:invariance}
  W = \mathrm{PI}^{\mathbb{B}}(p) + \mathrm{PI}^{\mathbb{A}}(p)
    = \sum_{k \in \mathbb{B}} \pi_k x_k
      - \sum_{k \in \mathbb{A}} \pi_k x_k ,
\end{equation}
which does not depend on $p$.

Expanding the two sums, the terms in $p$ enter as
$p\left(\sum_{\mathbb{A}} x_k - \sum_{\mathbb{B}} x_k\right)$, which
vanishes by flow conservation. Two consequences follow immediately.
First, $\partial \mathrm{PI}^{\mathbb{B}} / \partial p = -V^{*}$ and
$\partial \mathrm{PI}^{\mathbb{A}} / \partial p = +V^{*}$: every unit of
price improvement conferred on one side is subtracted from the other.
Second, a measured increase in buy-side price improvement is not evidence
that the mechanism created value: total gains from trade are fixed by the
allocation, and the clearing price only divides them. The metric must
therefore be read as a distributive statistic, reported for both sides
jointly. If $\mathcal{P}^{*}$ were a singleton, the division would be
pinned down by the submitted limits and there would be nothing to
contest; because it is generically an interval, the surplus split depends
on a rule, and any participant able to influence the position of
$\mathcal{P}^{*}$ or of $\sigma$ within it has an incentive to do so.

\subsection{The Last-Mover Advantage}

Suppose the engine operates under the fully open disclosure regime, so
that a participant submitting at time $t$ observes the schedules formed by
all orders with $t_k < t$.

\textbf{Late submission weakly dominates.} Suppose a participant can see
the orders already submitted by others, and hence the residual supply or
demand it faces, and that the selection rule $\sigma$ responds
monotonically to submitted limit prices, so that a participant can shift
the clearing price within $\mathcal{P}^{*}$ by adjusting its own limit.
Writing $\underline{p}=\min\mathcal{P}^{*}$ and
$\overline{p}=\max\mathcal{P}^{*}$, the participant's best response is to
set its limit so that the clearing price moves to the edge of
$\mathcal{P}^{*}$ most favourable to it,
\begin{equation}
\label{prop:lastmover}
  p^{*} \to \underline{p} \text{ for a buyer,} \qquad
  p^{*} \to \overline{p} \text{ for a seller,}
\end{equation}
while still keeping its order executed. Because this manoeuvre requires
knowing the other orders, a participant's surplus can only grow with the
information it has when it submits. Submitting later is therefore never
worse than submitting earlier, and it is strictly better whenever the
extra information changes the participant's optimal limit.

The argument is the standard one for sequential-move bargaining over a
fixed surplus, applied to the division established in
\Cref{lem:invariance}. A seller who observes a resting buy limit of
$\pi_{\mathbb{B}}$ knows that any ask weakly below it will execute, and
therefore submits just below $\pi_{\mathbb{B}}$ rather than at its own
valuation, capturing the interval. Truthfulness is not violated in its own
terms, since it is stated for a participant that cannot move the clearing
price; what disclosure does is make the boundaries of $\mathcal{P}^{*}$
computable, and hence make deviation actionable for any participant whose
order is large enough to reach them.

This identifies the source of the incentive to wait. When several prices
clear the same maximum volume, the choice among them is not neutral: it
transfers surplus between buyers and sellers. Under a monotone $\sigma$,
the participant who moves last can capture this transfer by pushing the
price to the boundary of $\mathcal{P}^{*}$ in its own favour. The prize
goes to whoever submits closest to the deadline with the fullest view of
the book. This is the \emph{boundary latency arbitrage}: a reward for
speed and timing rather than for providing liquidity or information.

Two conditions expose a participant to being the donating side of the
transfer: submitting without observing the accumulated batch, as an
uninformed participant does, and submitting early without repricing. Since
the transfer is zero-sum by the result above (\Cref{lem:invariance}), every participant exposed
in this way is matched by one with an incentive to exploit it. An order
resting in an open batch is thus a free option written to every subsequent
submitter in the window. This is the same option \citet{copeland1983information}
describe for a resting limit order in a continuous book: there, it is
written to traders reacting to news; here, the same option is written to
traders reacting to the accumulating batch itself.

\subsection{Cancellation, Repricing, and the Deadline Race}

The natural defence against the last-mover advantage above
(\Cref{prop:lastmover}) is to withdraw or reprice an order that has become
adversely positioned. This is where the speed race re-enters.

\textbf{Latency monotonicity.} Let participant $i$ face message latency
$\ell_i > 0$ and let the batch close deterministically at $\tau$. A
message dispatched at time $t$ is included in the batch if and only if
\begin{equation}
\label{prop:latency}
  t + \ell_i \le \tau,
\end{equation}
so the latest observable state on which $i$ can act is that prevailing at
$\tau - \ell_i$. The information set available to $i$ at the moment of
its final decision is therefore ordered by $\ell_i$, and expected surplus
is non-increasing in $\ell_i$.

The structure of the advantage differs from the continuous case but its
source does not. In a continuous book, lower latency wins a contest for
priority; in a batch, priority within the window is irrelevant by
construction, and lower latency instead buys a later decision point. The
fast trader observes more of the batch before committing, can reprice on
more recent information, and can place its final message closer to the
boundary while retaining the same probability of inclusion. Speed is not
neutralised; its payoff is relocated from ordering to deadline precision.
That some speed advantage survives batching is consistent with the
experimental evidence of \citet{aldrich2020experiments}, who find that
investment in speed falls but does not disappear under batch clearing, and
with the model of \citet{zoican2021speed}.

Three remarks bound the mechanism. Prohibiting cancellation within the
batch does not close the channel: a participant with two-sided access can
still neutralise an adverse position with an offsetting order, a
synthetic cancellation that needs no cancel message at all (true
enforcement would require account-level self-match prevention). And
banning cancellation outright would, if anything, strengthen the incentive
described above (\Cref{prop:lastmover}), by raising the cost of correcting
an early submission instead of removing it. The channel is also
disclosure-conditional: the informational component of the last-mover
advantage requires an observable accumulating batch and is absent under a
sealed-bid regime.

\subsection{Equilibrium Implication and Design Responses}

If late submission weakly dominates and lateness is bounded only by
$\ell_i$, submissions concentrate near the boundary and are ordered by
latency. In the limit the accumulation phase carries no orders and
therefore no information, and the auction degenerates into a sealed tender
conducted in the final microseconds of the window: the batch retains its
length nominally while losing the property the length was chosen to
deliver.

Three design responses follow. Randomising the close time makes waiting
costly, since an order dispatched at $t$ is included only if the realised
$\tau$ exceeds $t + \ell_i$; the participant then solves an interior
optimum trading the value of information against the risk of exclusion,
which justifies randomisation independently of the sniping argument
usually advanced for it. Restricting disclosure to an aggregate imbalance
or to nothing removes the informational component of the result above
(\Cref{prop:lastmover}) while leaving bid shading, so the choice is between
two costs rather than between one cost and none. Finally, and least
discussed in the literature, $\sigma$ can be made insensitive to the
marginal submitted limits, for instance by choosing the point of
$\mathcal{P}^{*}$ closest to an external reference price. This removes the
lever without altering the allocation, since by \Cref{lem:invariance} the
allocation is invariant on $\mathcal{P}^{*}$ in any case.
